\documentclass[11pt]{article}
\usepackage[utf8]{inputenc}
\usepackage{amsmath,amssymb}
\usepackage[margin=1in]{geometry}
\usepackage{hyperref}
\emergencystretch=3em

\title{Regrouping the merged expansion by exponent}
\author{Claude, with Daniel Murfet}
\date{2026-09-07}

\begin{document}
\maketitle
\sloppy

\begin{abstract}
The merged $d=2$ cutoff expansion, indexed by face indices $i<M_1$, $j<M_2$, is regrouped over the
finite pole set $\Lambda_M=\{\alpha_i\}_{i<M_1}\cup\{\delta_j\}_{j<M_2}$ as
$\sum_{\alpha\in\Lambda_M}N^{-\alpha}(A_\alpha\log N+B_\alpha)$ with the canonical
cutoff-independent coefficients of the previous unit (\texttt{mergedSum\_eq\_canon}). The
compatibility inequalities $(M_1-1)/k_1<M_2/k_2$ and $(M_2-1)/k_2<M_1/k_1$ are used exactly once
each: a retained pole of one axis lies strictly below every discarded pole of the other, so no
canonical coefficient is supported outside the retained ranges. The analytic cutoff theorems follow
in canonical pole form (\texttt{twoD\_cutoff\_canon}, \texttt{twoD\_cutoff\_amplitude\_canon}).
Zero \texttt{sorry}/\texttt{axiom}.
\end{abstract}

\section{The things formalised}
{\small
\begin{verbatim}
theorem uExp_lt_vExp_of (p : ℝ) (h₁ h₂ k₁ k₂ M₁ M₂ : ℕ) (hk₁ : 0 < k₁) (hk₂ : 0 < k₂)
    (hp₁ : ((h₁ : ℝ) + 1) / k₁ = p) (hp₂ : ((h₂ : ℝ) + 1) / k₂ = p)
    (hM₁ : ((M₁ : ℝ) - 1) / k₁ < (M₂ : ℝ) / k₂) (i j : ℕ) (hi : i < M₁) (hj : M₂ ≤ j) :
    uExp h₁ k₁ i < vExp h₂ k₂ j
theorem not_uPole_of_vExp ... (j : ℕ) (hj : j < M₂)
    (hno : ∀ i ∈ Finset.range M₁, uExp h₁ k₁ i ≠ vExp h₂ k₂ j) :
    ∀ i, uExp h₁ k₁ i ≠ vExp h₂ k₂ j
theorem uface_gamma_lt_of ... (i : ℕ) (hi : i < M₁) : (k₂ : ℝ) * uExp h₁ k₁ i - h₂ - 1 < M₂
noncomputable def poleSet (h₁ h₂ k₁ k₂ M₁ M₂ : ℕ) : Finset ℝ :=
  (Finset.range M₁).image (uExp h₁ k₁) ∪ (Finset.range M₂).image (vExp h₂ k₂)
theorem canonC_sum_eq ... :
    ∑ α ∈ poleSet h₁ h₂ k₁ k₂ M₁ M₂, transferTerm β α 1 (canonC h₁ h₂ k₁ k₂ c α) N
      = mergedLog β h₁ h₂ k₁ k₂ M₁ M₂ c N
theorem mergedSum_eq_canon ... (hU : ∀ i, i < M₁, canonical FP = FP at order M₂)
    (hV : ∀ j, j < M₂, canonical FP = FP at order M₁) (N : ℝ) :
    mergedSum β b h₁ h₂ k₁ k₂ M₁ M₂ a bj c N
      = ∑ α ∈ poleSet h₁ h₂ k₁ k₂ M₁ M₂,
          N ^ (-α) * (canonA β h₁ h₂ k₁ k₂ c α * Real.log N
            + canonB β b h₁ h₂ k₁ k₂ a bj c α)
theorem twoD_cutoff_amplitude_canon (β b p ρ : ℝ) (h₁ h₂ k₁ k₂ M₁ M₂ : ℕ) (hβ : 0 < β)
    (hb : 0 < b) (hbρ : b < ρ) (hk₁ : 0 < k₁) (hk₂ : 0 < k₂) (hp₁ : ((h₁ : ℝ) + 1) / k₁ = p)
    (hp₂ : ((h₂ : ℝ) + 1) / k₂ = p)
    (hM₁ : ((M₁ : ℝ) - 1) / k₁ < (M₂ : ℝ) / k₂) (hM₂ : ((M₂ : ℝ) - 1) / k₂ < (M₁ : ℝ) / k₁)
    (x y : ℕ × ℕ → ℝ) (hx : WSummable ρ x) (hy : WSummable ρ y) :
    (fun N : ℝ => twoDAmp β b N h₁ h₂ k₁ k₂ (anaAmp (ampCoeff β x y) b)
        - ∑ α ∈ poleSet h₁ h₂ k₁ k₂ M₁ M₂,
            N ^ (-α) * (canonA β h₁ h₂ k₁ k₂ (fun i j s => ampCoeff β x y (i, j) s) α
                * Real.log N
              + canonB β b h₁ h₂ k₁ k₂ (anaFaceU (ampCoeff β x y) b)
                  (anaFaceV (ampCoeff β x y) b) (fun i j s => ampCoeff β x y (i, j) s) α))
      =O[atTop] fun N : ℝ =>
        N ^ (-(p + min ((M₁ : ℝ) / k₁) ((M₂ : ℝ) / k₂))) * (1 + Real.log N)
\end{verbatim}
}

\section{Mathematics}

Each canonical pole term splits into the $u$-face, $v$-face and collision contributions
(\texttt{canon\_term\_eq}). The $u$-face block over $\Lambda_M$ reduces to the block over the
$u$-poles because $U_\alpha$ vanishes at a $v$-pole $\delta_j$ ($j<M_2$) which is no retained
$u$-pole: it is then no $u$-pole at all, since $\alpha_i\ge\alpha_{M_1}=p+M_1/k_1>p+(M_2-1)/k_2\ge\delta_j$
for $i\ge M_1$. Over the $u$-poles the sum pulls back along the injective map $i\mapsto\alpha_i$
and the canonical finite part is replaced by the order-$M_2$ one (admissible since
$\gamma_i=k_2\alpha_i-h_2-1<M_2$, again by compatibility). The collision block is the merged log
family: for each $i<M_1$ either a unique $j_0<M_2$ collides (then the indicator sum is that single
term and $C_{\alpha_i}=c_{ij_0}/(k_1k_2)$) or none does, in which case $\alpha_i$ is no $v$-pole at
all (compatibility once more) and both sides vanish.

\section{Paper-facing meaning}

With this the analytic $d=2$ theorem reads exactly like \texttt{thm:TaylorTree} restricted to a
finite pole set: $Z(N)=\sum_{\alpha\in\Lambda_M}N^{-\alpha}P_\alpha(\log N)+O(N^{-q_M}(1+\log N))$
with $P_\alpha$ of degree $\le1$ and coefficients defined once for all. The remaining step to the
paper's ``for every $T$'' form is the choice $M_\ell=\lceil k_\ell R\rceil$, $R=\max(1,2T-p)$, which
satisfies both compatibility inequalities and $q_M\ge2T$, plus discarding the retained poles
$\ge2T$ into the remainder.

\section{Lean notes}

\begin{itemize}
\item \texttt{Finset.sum\_subset} composed as
\texttt{((Finset.sum\_subset Finset.subset\_union\_left ?\_).symm).trans ?\_} turns a sum over a
union into a sum over one part plus a vanishing obligation; then \texttt{Finset.sum\_image} with the
axis injectivity pulls back to \texttt{range}.
\item The \texttt{if} conditions in \texttt{mergedLog} are spelled with the literal exponent
expressions; a \texttt{change} to the \texttt{uExp}/\texttt{vExp} form (definitional) is needed
before \texttt{if\_pos}/\texttt{if\_neg} with hypotheses stated via \texttt{uExp}.
\item \texttt{gcongr} proves $i/k_1\le(M_1-1)/k_1$ from the side goal $i\le M_1-1$ (discharged by
\texttt{linarith} from the cast of $i+1\le M_1$).
\item \texttt{lt\_div\_iff₀} is the current name for clearing a positive denominator in a strict
inequality.
\item The unit compiled on the first attempt.
\end{itemize}

\end{document}
